\section{Context}

\subsection{Curiosity Origin}
As frequent travelers, we are often curious about how people move through large cities.
In particular, we wanted to understand mobility patterns, the choices people make when navigating urban spaces, and the factors that shape those choices.

This curiosity led us to explore data sources that could provide meaningful insights into urban movement.
Among them, New York City's bike rental data stood out as a rich source of information on how people travel across the city using a sustainable mode of transportation.
When combined with other datasets, it can offer a clearer picture of the city's daily dynamics.

We also learned that this bike rental data is provided by a private company (\texttt{Citi Bike}) that operates the bike-sharing service in New York City.
For this reason, making our analysis useful and relevant to the company is an effective way to increase its practical impact.

\subsection{Project Circumstance and Purpose}
\texttt{Citi Bike} currently collects vast amounts of operational data. 
If effectively analyzed and visualized, this data could provide valuable insights to optimize their services and enhance user experience.

Our visualization tool directly addresses these gaps by making complex mobility data accessible and actionable.
Our primary objective is to transform raw \texttt{Citi Bike} operational data into clear, actionable insights for stakeholders.
Specifically, we aim to reveal:\begin{itemize}
  \item \textbf{Temporal patterns} -- peak usage hours, seasonal trends, and demand cyclicity
  \item \textbf{Spatial hotspots} -- high-demand stations, commute corridors, and underutilized infrastructure
  \item \textbf{User segmentation} -- behavioral differences between members and casual riders
\end{itemize}

\subsubsection*{Target Audience}
Our project is primarily designed for the \texttt{Citi Bike} operations team.
Although, we intend to also target secondary audiences, such as city planners and the general public.

The idea is to create a tool that provides for each of these audiences the following benefits:
\begin{itemize}
  \item \textbf{Citi Bike operations} -- quickly identify optimization opportunities for station placement and bike rebalancing.
  This allows them to do data-driven decisions to improve the efficiency of their service, such as optimizing bike distribution across stations or identifying underperforming locations.
  \item \textbf{City planners} -- understand commute flows and modal patterns to inform infrastructure investment and urban planning decisions.
  \item \textbf{General public} -- develop awareness of urban mobility and sustainable transportation alternatives
\end{itemize}

\subsubsection*{Constraints}
%#TODO  Pressure: how much time do you have and what milestones along the realisation path?
%       rules: layout/size restrictions, style guidelines, functional restrictions

\subsubsection*{Deliverables Vision}
We deliver an interactive web-based dashboard as the primary output.
This allows all users to explore mobility patterns through a unified, user-friendly interface.
Our design approach is to present clear, immediate insights at first impactful while also allowing deeper exploration for users who want to examine the details.

Regarding update frequency, we expect the dashboard to be refreshed automatically on a regular basis, since the data are continuously collected.
This ensures that the insights remain useful and relevant over time, while also enabling historical analysis and comparisons across different periods.

\subsubsection*{Design Philosophy}
Our visualization prioritizes clarity and analytical depth over visual polish.
We employ a dual-level approach: high-level dashboards present immediate summary insights (e.g., daily ridership totals, top stations), while detailed views enable granular exploration (e.g., hourly breakdowns, user type comparisons).
All design choices remain data-centric, avoiding decorative elements or misleading visual encodings that could obscure underlying patterns.

\subsubsection*{Resources}
Since we are working with a large dataset and aim to create an interactive dashboard, we need programming skills and experience with data and visualization tools.

Our tech stack combines:
\begin{itemize}
  \item \textbf{Backend (Python)} -- efficient data processing, cleaning, and REST API development via FastAPI;
  \item \textbf{Frontend (JavaScript/D3.js)} -- custom, interactive visualizations tailored to mobility pattern analysis.
\end{itemize}
This modular architecture enables rapid iteration during development while remaining scalable as data volumes grow.

